\setauthor{Sebastian Radić \& Luis Schörgendorfer}
\chapter{Schlussbetrachtung}
\label{chap:schlussbetrachtung}

\section{Zusammenfassung der Ergebnisse und Beiträge}
Das Ziel dieser Diplomarbeit war die Entwicklung eines Systems, das eingehende
Rechnungen -- egal ob als PDF oder als Bild -- automatisch in maschinenlesbare
XML-Formate umwandelt. Dieser Prototyp liegt jetzt vor, und er funktioniert.

Alle zehn Testrechnungen, die die PROGRAMMIERFABRIK GmbH für die Abnahme
bereitgestellt hat, wurden vom \textit{SmartBillConverter} vollständig verarbeitet.
Die generierten XML-Ausgaben haben sowohl den offiziellen ebInterface-Validator
der AUSTRIAPRO als auch den deutschen E-Rechnungs-Validator für ZUGFeRD~2.3
fehlerfrei bestanden. Das war das definierte Erfolgskriterium -- und es wurde erreicht.

Technisch setzt das System auf eine durchgängige Verarbeitungspipeline:
Textextraktion aus PDF (PdfPig) oder Bild (Tesseract~OCR), anschließende
Normalisierung durch die Google Gemini~API in eine strukturierte
\texttt{InvoiceData}-Darstellung und abschließende XML-Generierung im
jeweiligen Zielformat. Das Angular-Frontend bietet eine einfach bedienbare
Oberfläche, die den gesamten Ablauf abbildet -- vom Upload bis zum Download
des fertigen~XMLs.

Sebastian Radić hat das Backend entworfen und implementiert: die gesamte
Verarbeitungspipeline, die Service- und Repository-Schicht, die Datenbankanbindung
über Entity Framework Core mit PostgreSQL sowie die theoretischen Grundlagen
zu den E-Rechnungsstandards. Luis Schörgendorfer hat das Frontend entwickelt:
Anforderungsanalyse, Angular-Architektur, alle UI-Komponenten, die Backend-Anbindung
und die Evaluation der fertigen Anwendung. Die gemeinsamen Kapitel -- Einleitung,
Gesamtarchitektur, Technologiestack und Schlussbetrachtung -- wurden von
beiden Autoren erarbeitet.

\section{Ausblick und Übertragbarkeit}
Der \textit{SmartBillConverter} wurde als internes Tool der PROGRAMMIERFABRIK GmbH
entwickelt, aber der zugrundeliegende Ansatz ist nicht unternehmensspezifisch.
Die Pipeline aus OCR, KI-Normalisierung und schema-gesteuerter XML-Generierung
lässt sich auf jedes strukturierte Dokumentformat übertragen, bei dem aus
unstrukturierten Eingaben ein definiertes Ausgabeformat entstehen soll.

Die praktische Erfahrung zeigt, dass Sprachmodelle gut darin sind, inhaltliche
Zusammenhänge in Freitext zu erkennen und in strukturierte Daten zu überführen --
auch dann, wenn das Eingabedokument keinerlei einheitliches Layout hat.
Gleichzeitig bleibt das Ergebnis fehleranfällig genug, dass eine manuelle
Überprüfung sinnvoll bleibt. Der Mensch bleibt vorerst im Loop.

Mit den Erweiterungen, die im vorangegangenen Kapitel beschrieben wurden --
lokales KI-Modell, weitere Standards, Batch-Import und Feedback-Loop --
könnte aus dem Prototyp ein ernsthafter Produktivbetriebskandidat werden.
Die Grundlage dafür ist gelegt.
